\scanpage{126}
s'écrit
\[
y_1y'_1+y_2y'_2=1
\]
et on a
\[
t=y_1^{\deleted{\unsure{'|1}}}y'_1
\]
\deleted{On peut prendre} les relations deviennent \deleted{$y_1=1,y'_1=t$ et $y_2,y'_2$ tel que $y_1y'_1=1-t$},
\[
\left\{\begin{aligned}y_1y'_1&=t\\ y_2y'_2&=1-t\end{aligned}\right.
\]

Si $t\ne0,1$ on peut choisir arbitrairement $y_1,y_2$ \deleted{$\unsure{x|\in}$} $\in k^{\times}$, et tirer \deleted{$y'_1,\unsure{y'|y'_2}$}
\[
y'_1=t/y_1,\quad y'_2=\frac{1-t}{y'_2}
\]
Si $t=0$, \deleted{on a $y_1$ \unsure{.}} ça devient
\[
\begin{aligned} y_1y'_1&=0\\ y_2y'_2&=1\end{aligned}
\]
deux familles de solutions \inserted{à deux par.} suivant que $y_1=0$ ou $y'_1=0$ et symétriques si $t=1$.

\uline{Cas 2} On peut tjrs $x=e_1$, $x'=e'_1$, \unsure{relation|relations}
\[
\begin{aligned} y_1y'_1+y_2y'_2&=0\\ t&=y_1y'_1\end{aligned}
\]
i.e. comme équations en les $y_i,y'_i$ : \deleted{$\unsure{y_2}$}
\[
\left\{\begin{aligned} y_1y'_1&=t\\ y_2\unsure{y'_2|y_2^2}&=-t\end{aligned}\right.
\]

\uline{Cas 3} On peut prendre $x=e_1$, $x'=e'_{\underline{\underline 2}}$.
\[
\deleted{\unsure{y_1y'_1+y_2y'_2=0}}
\]
On doit avoir
\[
\left\{\begin{aligned} y_1y'_1+y_2y'_2&=0\\ t&=y'_1y_2\end{aligned}\right.
\]
\deleted{\unsure{et}} \ldots
% archive sheet number: 125
